\section{Conclusion}
\label{sec:conclusion}

We argued that the standing triage decision in agentic memory --- what to
encode, forget, and retrieve under a fixed budget --- is mis-served by
the single-factor similarity and recency policies in common use, because
the forgetting decision is made before the future query is known. We
proposed a multi-factor memory value $V(m)=\sum_i w_i f_i(m)$ whose
weights are learned from downstream task return and whose single scalar
uniformly governs all three operations. On the real LongMemEval
benchmark, a learned multi-factor value retains $0.811\pm0.047$ of gold
evidence in the realistic blind-forgetting regime, versus $0.633$ for
uniform weighting, $0.380$ for recency, and at most $0.39$ for any single
factor, with interpretable weights that recover the workload's true
signal. We also showed why this advantage is invisible to a query-defined
retention metric: under an oracle that peeks at the evaluation question,
similarity alone saturates retention and the multi-factor gap vanishes
--- a distinction any honest evaluation of forgetting must draw.

The immediate next step is to close the loop to end-task accuracy:
fitting the weights to downstream QA correctness (unlocking the
task-utility factor) and measuring answer quality, not just retention,
under each policy. Beyond that, the same learned-value harness applied to
tool-use and coding agents would test whether the value-directed view of
memory transfers across agentic workloads. The substrate, factor
annotator, and evaluation harness are open-source, and every number in
this paper reproduces on a single CPU.
